\section*{第 \ref{ch:b1:vspaces} 章 --- 向量空间}

\begin{solution}{exo:b1:vspaces:1}
\begin{enumerate}
  \item 是：含 $0$，且定义方程是线性的
        （对 $x + \lambda y$ 保持稳定）。
  \item 否：不含 $(0,0,0)$。
  \item 否：$(1, 0)$ 与 $(0, -1)$ 属于（$xy = 0$），它们的和 $(1,
        -1)$ 不属于（$xy = -1 < 0$）。
  \item 是：$0$ 在 $1$ 处为零；$(P + \lambda Q)(1) = P(1) +
        \lambda Q(1) = 0$。
  \item 是：零函数有界；若 $\abs f \leq M$ 且
        $\abs g \leq M'$，则 $\abs{f + \lambda g} \leq M +
        \abs\lambda M'$。
\end{enumerate}
\end{solution}

\begin{solution}{exo:b1:vspaces:2}
$(1,2,1) = a(1,0,1) + b(1,1,0)$ 要求 $a + b = 1$，$b = 2$，$a =
1$：不相容（$a + b = 3 \neq 1$），故不在其张成的子空间中。$(2,1,1) =
a(1,0,1) + b(1,1,0)$：$b = 1$，$a = 1$，$a + b = 2$，相容，因此
$(2,1,1) = (1,0,1) + (1,1,0)$，属于该张成的子空间。

方程：由 $(x, y, z) = (a + b, b, a)$ 得 $x = y + z$，故张成的子空间为
平面 $\{x - y - z = 0\}$。
\end{solution}

\begin{solution}{exo:b1:vspaces:3}
第一组向量：由 $\lambda(1,1,0) + \mu(1,0,1) + \nu(0,1,1) = 0$ 得
$\lambda + \mu = 0$，$\lambda + \nu = 0$，$\mu + \nu = 0$，相加得
$2(\lambda + \mu + \nu) = 0$，再将各式逐一相减，
$\lambda = \mu = \nu = 0$，故线性无关。

第二组：$(2,4,6) = 2(1,2,3)$，线性相关。

第三组：$\R^3$ 中的四个向量 --- 由维数可知（见
\cref{ch:b1:findim}）必然线性相关；直接验证：$(0,1,1) = -(1,0,0) +
0\cdot(1,1,0) + (1,1,1)$，事实上 $(-1,0,0) + (1,1,1) = (0,1,1)$，这是
一个非平凡的线性关系。
\end{solution}

\begin{solution}{exo:b1:vspaces:4}
多项式 $1, (X-1), (X-1)^2, (X-1)^3$ 的次数互不相同，分别为 $0,
1, 2, 3$，故线性无关（由\cref{prop:b1:vspaces:freecriteria}~(1)可知）；四个线性无关的
向量张成整个空间（每个 $P \in \R_3[X]$ 都可按 $X -
1$ 的幂展开，例如由多项式的泰勒公式，参见
\cref{prop:b1:poly:multiplicity}），构成一组基。对 $X^3$，取其在
$1$ 处的泰勒展开：$P = X^3$，$P(1) = 1$，$P'(1) = 3$，$P''(1) = 6$，$P'''(1) =
6$：
\[
X^3 = 1 + 3(X - 1) + 3(X-1)^2 + (X-1)^3 ,
\]
坐标为 $(1, 3, 3, 1)$（帕斯卡三角的一行，与 $X^3 =
((X-1)+1)^3$ 相符）。
\end{solution}

\begin{solution}{exo:b1:vspaces:5}
$F \cap G$：条件 $x = y = z$ 与 $x = t = 0$ 合起来得
$x = 0$，从而 $y = z = 0$，且 $t = 0$，故交集为
$\{0\}$。和：对给定的 $(x,y,z,t)$，求 $(a,a,a,b) \in F$ 与
$(0,c,d,0) \in G$ 使其和为该向量：$a = x$，$b = t$，$c = y - x$，$d =
z - x$，总是可以取到。故 $\R^4 = F \oplus G$，且
\[
(1,2,3,4) = (1,1,1,4) + (0,1,2,0) .
\]
\end{solution}

\begin{solution}{exo:b1:vspaces:6}
$G$ 中的元素为 $\lambda u$；若其收敛，则（因
$\lambda u_n = \lambda(-1)^n$ 有两个子列极限 $\pm
\lambda$），必有 $\lambda = 0$，即 $F \cap G = \{0\}$。

$F + G$ 并非全部：它由形如
$c_n + \lambda(-1)^n$ 的数列组成，其中 $(c_n)$ 收敛。数列 $v_n =
n$ 不具有此形式（$v_n - \lambda(-1)^n$ 无界，从不
收敛）。因此 $F \oplus G \subsetneq$（所有数列构成的空间）。
\end{solution}

\begin{solution}{exo:b1:vspaces:7}
（$\supseteq$）$F \cap G$ 与 $F \cap H$ 均含于 $F$，且它们的
和含于 $G + (F \cap H)$，故包含关系成立，因为左端
是一个包含上述两部分的子空间 --- 具体地，任取元素 $g +
h$，其中 $g \in F\cap G$，$h \in F \cap H$，它属于 $F$（两个
$F$ 中元素之和），也属于 $G + (F \cap H)$。

（$\subseteq$）设 $x \in F$，且 $x = g + h$，其中 $g \in G$，$h \in F
\cap H$。于是 $g = x - h \in F$（$F$ 中元素之差），从而
$g \in F \cap G$，且 $x = g + h \in (F \cap G) + (F \cap H)$。

完全分配律在 $\R^2$ 中的反例：$F =
\operatorname{Vect}(1,1)$, $G = \operatorname{Vect}(1,0)$, $H =
\operatorname{Vect}(0,1)$。则 $G + H = \R^2$，故 $F \cap (G+H) =
F$，而 $F \cap G = F \cap H = \{0\}$，右端为 $\{0\}
\neq F$。
\end{solution}

\begin{solution}{exo:b1:vspaces:8}
\begin{enumerate}
  \item 设 $\sum_{i} \lambda_i \eu^{a_i x} = 0$ 对所有 $x$ 成立。
        乘以 $\eu^{-a_p x}$，得 $\lambda_p + \sum_{i < p}
        \lambda_i \eu^{(a_i - a_p)x} \to \lambda_p$，当 $x \to
        +\infty$ 时（每个指数 $a_i - a_p < 0$）。左端
        恒为 $0$，故 $\lambda_p = 0$；依此向下递推。
  \item 设 $a\cos x + b \sin x + c \cos 2x + d\sin 2x = 0$ 对所有
        $x$ 成立。取 $x = 0$，得 $a + c = 0$；取 $x = \pi$，得 $-a + c
        = 0$；从而 $a = c = 0$，关系式化为 $b\sin
        x + d \sin 2x = 0$。取 $x = \frac\pi2$，得 $b = 0$；
        再取 $x = \frac\pi4$，得 $d = 0$。
  \item 设 $\sum \lambda_i \abs{x - a_i} = 0$ 对所有 $x$ 成立。
        函数 $\sum_{i \neq j} \lambda_i\abs{x - a_i}$ 在
        $a_j$ 处可导（除其自身角点外，每一项均可导），
        故 $-\lambda_j \abs{x - a_j}$（二者之差）
        必须在 $a_j$ 处也可导 --- 这迫使
        $\lambda_j = 0$（$\abs{\,\cdot\,}$ 有角点）。此式对
        每个 $j$ 均成立。
\end{enumerate}
\end{solution}

\begin{solution}{exo:b1:vspaces:9}
设 $h \in H$。由 $h \in H \subseteq F + H = F + G$，可写 $h = f
+ g$，其中 $f \in F$，$g \in G$。于是 $f = h - g \in H$（两项均在
$H$ 中，用到 $G \subseteq H$），于是 $f \in F \cap H = F \cap G
\subseteq G$，且 $h = f + g \in G$。因此 $H \subseteq G$，结合
假设 $G \subseteq H$，二者相等。

反例（不假设 $G \subseteq H$）：在 $\R^2$ 中取 $F =
\operatorname{Vect}(1,0)$，$G = \operatorname{Vect}(0,1)$，$H =
\operatorname{Vect}(1,1)$，则 $F + G = F + H = \R^2$，且 $F \cap G
= F \cap H = \{0\}$，但 $G \neq H$。
\end{solution}

\begin{solution}{exo:b1:vspaces:10}
两个集合都是子空间（定义条件是线性的，
且对 $0$ 成立）。分解：对 $P \in \R[X]$，
\[
P(X) = \underbrace{\frac{P(X) + P(-X)}{2}}_{\in\,\mathcal P}
+ \underbrace{\frac{P(X) - P(-X)}{2}}_{\in\,\mathcal I},
\]
且既偶又奇的多项式满足 $P = -P$，故 $P = 0$：
因此和是直和，且等于 $\R[X]$。

现设 $P = \sum_k a_k X^k$ 为偶多项式。则 $P(X) - P(-X) = 2
\sum_{k \text{ 奇数}} a_k X^k$ 为零多项式，故每个
奇次项系数均为零（见\cref{def:b1:poly:def}）：$P \in
\operatorname{Vect}(1, X^2, X^4, \dots)$，即 $P = Q(X^2)$，其中 $Q$ 为
某个多项式。反之，任一关于 $X^2$ 的多项式都是偶多项式。
\end{solution}

\begin{solution}{exo:b1:vspaces:11}
设 $a(x_1 + x_2) + b(x_2 + x_3) + c(x_3 + x_1) = 0$。
对线性无关族 $(x_1, x_2, x_3)$ 重新合并：
\[
(a + c)\,x_1 + (a + b)\,x_2 + (b + c)\,x_3 = 0
\implies a + c = a + b = b + c = 0 .
\]
前两式相减得 $c = b$；再由第三式
得 $2b = 0$，故 $b = c = 0$，进而 $a = 0$，因此该族线性无关。

对四个向量，类似的族\emph{总是}线性相关：
\[
(x_1 + x_2) - (x_2 + x_3) + (x_3 + x_4) - (x_4 + x_1) = 0
\]
是一个非平凡的零组合（系数为 $1, -1, 1, -1$），
无论 $(x_1, x_2, x_3, x_4)$ 取何值。圈长的奇偶性
起决定作用。
\end{solution}

\begin{solution}{exo:b1:vspaces:12}
\begin{enumerate}
  \item 若 $F_1 \subseteq F_2$ 或 $F_2 \subseteq F_1$，则并集
        是其中之一，从而是真子集。否则取 $x \in F_1
        \setminus F_2$ 及 $y \in F_2 \setminus F_1$，考虑
        $x + y$。若 $x + y \in F_1$，则 $y = (x + y) - x \in
        F_1$，矛盾。若 $x + y \in F_2$，则 $x \in F_2$：
        矛盾。因此 $x + y \notin F_1 \cup F_2$，且 $E \neq
        F_1 \cup F_2$。
  \item 用反证法，假设 $E = F_1 \cup \dots \cup
        F_k$，取 $k$ 为所有此类覆盖中\emph{最小}者。
        最小性排除 $F_1 \subseteq F_2 \cup
        \dots \cup F_k$（否则可去掉 $F_1$），故存在 $x \in F_1$
        使 $x \notin F_i$ 对所有 $i \geq 2$ 成立。由于 $F_1$ 为
        真子空间，取 $y \notin F_1$。对每个标量 $t$，
        向量 $y + t x$ 落在某个 $F_i$ 中。它绝不落在
        $F_1$ 中：否则 $y = (y + tx) - tx \in F_1$（因 $x \in
        F_1$）。域是无限的，故可取 $k$ 个互异的
        标量 $t_1, \dots, t_k$：这 $k$ 个向量 $y + t_j x$
        落入 $k - 1$ 个子空间 $F_2, \dots, F_k$ 中，于是有两个
        其中两个不妨设为 $y + t x$ 与 $y + t' x$（$t \neq t'$），
        落在同一个 $F_i$ 中（$i \geq 2$）。于是它们的差
        $(t - t')x \in F_i$，故 $x \in F_i$，矛盾。因此
        不存在由有限个真子空间构成的覆盖。
\end{enumerate}
\end{solution}

\begin{solution}{pb:b1:vspaces:1}
\textbf{1.} $L_i$ 是 $n$ 个一次因式之积除以一个
非零常数（$x_i$ 两两不同），故 $\deg L_i = n$。
在 $x_j$（$j \neq i$）处求值：分子中的因式 $X - x_j$
为零，故 $L_i(x_j) = 0$。在 $x_i$ 处，分子与分母
恰好相同，故 $L_i(x_i) = 1$。

\textbf{2.} 设 $\sum_i \lambda_i L_i = 0$。在 $x_j$ 处求值：
除 $\lambda_j L_j(x_j) = \lambda_j$ 外，其余各项均为零，故
$\lambda_j = 0$（对每个 $j$），即该族线性无关。

\textbf{3.} 令 $D = P - \sum_i P(x_i) L_i$。则 $\deg D \leq n$，且由第 1 题，
$D(x_j) = P(x_j) - P(x_j) = 0$，对 $n + 1$
个互异点 $x_0, \dots, x_n$ 成立。次数
$\leq n$ 的非零多项式至多有 $n$ 个根（见\cref{cor:b1:poly:nroots}），故 $D
= 0$。因此每个 $P \in \R_n[X]$ 都是 $L_i$ 的线性组合：
即该族为生成族，结合第 2 题，它是一组基。

\textbf{4.} 给定 $y_0, \dots, y_n$，多项式 $P = \sum_i
y_i L_i$ 的次数 $\leq n$，且满足插值条件。唯一性：任一
插值多项式 $P$ 由第 3 题，坐标为 $(P(x_0), \dots,
P(x_n)) = (y_0, \dots, y_n)$（基为 $(L_i)$），而
基下坐标唯一（见\cref{prop:b1:vspaces:coordinates}）。$P$
在拉格朗日基下的坐标，正是它在\emph{节点处的
取值}——这正是该基的意义所在。

\textbf{5.} 对 $P = X^k$（$0 \leq k \leq n$）应用第 3 题：
\[
X^k = \sum_{i=0}^{n} x_i^{k} L_i ,
\]
而 $k = 0$ 时得 $\sum_i L_i = 1$。

\textbf{6.} $\deg N_k = k$ 恰成立：族 $(N_0, \dots, N_n)$
是 $\R_n[X]$ 中次数呈阶梯状的族，由
\cref{ex:b1:vspaces:staircase}可知它是基（线性无关性由
\cref{prop:b1:vspaces:freecriteria}~(1) 得出，生成性由对次数的有限
下降法得出）。

\textbf{7.} 由递推关系，
\[
f[x_0, x_1] = \frac{f(x_1) - f(x_0)}{x_1 - x_0},
\qquad
f[x_0, x_1, x_2] = \frac{f[x_1, x_2] - f[x_0, x_1]}{x_2 - x_0}.
\]
对 $f(x) = x^2$：
\[
f[x_0, x_1] = \frac{x_1^2 - x_0^2}{x_1 - x_0} = x_0 + x_1,
\]
于是
\[
f[x_0, x_1, x_2]
= \frac{(x_1 + x_2) - (x_0 + x_1)}{x_2 - x_0}
= \frac{x_2 - x_0}{x_2 - x_0} = 1 .
\]

\textbf{8.} $\deg S \leq n$，因为 $Q, R$ 的次数均 $\leq n - 1$。
在 $x_0$ 处：$S(x_0) = \frac{-(x_0 - x_n) R(x_0)}{x_n - x_0} =
R(x_0) = f(x_0)$。在 $x_n$ 处：$S(x_n) = \frac{(x_n - x_0)
Q(x_n)}{x_n - x_0} = Q(x_n) = f(x_n)$。在内部结点 $x_i$
（$1 \leq i \leq n-1$），$Q$ 与 $R$ 的取值均为 $f(x_i)$，
于是
\[
S(x_i) = \frac{(x_i - x_0) - (x_i - x_n)}{x_n - x_0}\, f(x_i)
= f(x_i) .
\]

\textbf{9.} 对结点个数作归纳。一个结点时，插值多项式
即常数 $f(x_0) = f[x_0]$。假定结论对
$k$ 个结点成立，设 $S$ 在 $x_0, \dots, x_k$ 处插值。由唯一性
（第 4 题），$S$ 可由艾特肯引理从 $R$
（结点 $x_0, \dots, x_{k-1}$）与 $Q$（结点 $x_1, \dots, x_k$）给出。
$X^{k}$ 在 $S$ 中的系数为
\[
\frac{[X^{k-1}]\,Q - [X^{k-1}]\,R}{x_k - x_0}
= \frac{f[x_1, \dots, x_k] - f[x_0, \dots, x_{k-1}]}{x_k - x_0}
= f[x_0, \dots, x_k]
\]
由归纳假设及定义递推关系可得。

\textbf{10.} 设 $P_k$ 为 $f$ 在 $x_0, \dots, x_k$ 处的插值多项式。
该差 $P_k - P_{k-1}$ 的次数 $\leq k$，且在
$x_0, \dots, x_{k-1}$ 处为零，故由因式定理连用 $k$ 次
（\cref{thm:b1:poly:factor}）得其等于 $c\,N_k$，$c$ 为常数；
比较 $X^{k}$ 的系数并利用第 9 题，$c =
f[x_0, \dots, x_k]$。由 $P_0 = f(x_0) N_0$ 逐项相消，得
牛顿公式。求显式表达式时，记 $P_k = \sum_{i \leq
k} f(x_i) L_i$（拉格朗日形式，结点 $x_0, \dots, x_k$），并读取
$X^{k}$ 的系数：每个 $L_i$ 的贡献为
$\frac{1}{\prod_{j \neq i}(x_i - x_j)}$，于是
\[
f[x_0, \dots, x_k] = \sum_{i=0}^{k}
\frac{f(x_i)}{\prod_{j \neq i,\, j \leq k}(x_i - x_j)} .
\]
右端在结点的任意置换下不变，故
差商与结点的排列次序无关。

\textbf{11.} 若 $P = a X^m + (\text{低次项})$，则由二项式
定理得
\[
\Delta P = a\bigl((X+1)^m - X^m\bigr) + \dots
= a\,m\,X^{m-1} + (\text{低次项}),
\]
因为 $(X+1)^m - X^m = m X^{m-1} + \dots$，且
$P$ 的低次部分，其次数 $\leq m - 2$，经 $\Delta$ 后（或
次数 $\leq m-2$ 的项）。故 $\deg \Delta P = m - 1$，首项
系数为 $m a$。常数 $c$ 给出 $\Delta c = c - c = 0$。

\textbf{12.} $\deg B_k = k$，为阶梯式，因而是
$\R_n[X]$ 的一组基（\cref{ex:b1:vspaces:staircase}）。对 $\Delta B_k$
（$k \geq 1$），提出公因式：
\begin{align*}
k!\,\Delta B_k
&= (X+1)X\cdots(X-k+2) - X(X-1)\cdots(X-k+1) \\
&= X(X-1)\cdots(X-k+2)\,\bigl[(X+1) - (X-k+1)\bigr] \\
&= k\,X(X-1)\cdots(X-k+2),
\end{align*}
于是 $\Delta B_k = \frac{X(X-1)\cdots(X-k+2)}{(k-1)!} = B_{k-1}$。

\textbf{13.} 记 $P = \sum_{k=0}^{n} c_k B_k$（基，第
12 题）。作用 $\Delta^{j}$：由第 12 题，$\Delta^{j} P = \sum_{k
\geq j} c_k B_{k-j}$。在 $0$ 处取值：$B_0(0) = 1$，且 $B_m(0) =
0$（$m \geq 1$，因子 $X$ 使其为零），于是
$\bigl(\Delta^{j}P\bigr)(0) = c_j$。这正是向前
差分公式。

\textbf{14.} 对 $k$ 作归纳。$k = 0$ 时，等式即为
$P(0) = P(0)$。假定对 $k$ 成立，将其用于 $\Delta P$：
\[
\bigl(\Delta^{k+1} P\bigr)(0)
= \sum_{j=0}^{k} (-1)^{k-j}\binom kj \bigl(P(j+1) - P(j)\bigr).
\]
合并 $P(i)$ 的系数：其为 $(-1)^{k-i+1}\binom
k{i-1}\cdot(-1)^{0}$，来自第一个和式（经平移），以及
$-(-1)^{k-i}\binom ki$ 来自第二个和式 --- 合起来
\[
(-1)^{k+1-i}\Bigl(\binom k{i-1} + \binom ki\Bigr)
= (-1)^{k+1-i}\binom{k+1}i
\]
由帕斯卡法则，这正是 $k + 1$ 时的恒等式。

\textbf{15.} 从次数 $n$ 出发反复应用第 11 题，首项系数为
$a_n$：作一次 $\Delta$ 后次数为 $n-1$，首项系数为
$n a_n$；作两次后为 $n(n-1)a_n$；作 $n$ 次之后，
次数为 $0$，值为 $n(n-1)\cdots 1\, a_n = n!\,a_n$，是一个常数。
再作一次 $\Delta$，常数即被消去：$\Delta^{n+1}P = 0$。

\textbf{16.} 若 $m \geq k$，则 $B_k(m) = \binom mk \in \N$。若 $0
\leq m < k$，则 $m(m-1)\cdots(m-k+1)$ 中有一个因子为零，故
$B_k(m) = 0$。若 $m = -q$，其中 $q \geq 1$：
\[
B_k(-q) = \frac{(-q)(-q-1)\cdots(-q-k+1)}{k!}
= (-1)^k\,\frac{q(q+1)\cdots(q+k-1)}{k!}
= (-1)^k \binom{q+k-1}{k},
\]
为整数。因此每个 $B_k$ 都把 $\Z$ 映入 $\Z$。

\textbf{17.} （$\Leftarrow$）若 $P = \sum_k c_k B_k$，且 $c_k \in
\Z$，则对 $m \in \Z$，$P(m) = \sum_k c_k B_k(m) \in \Z$，由
第 16 题。（$\Rightarrow$）若 $P$ 为整值多项式，则其
坐标为 $c_k = \bigl(\Delta^k P\bigr)(0) = \sum_{j=0}^k
(-1)^{k-j}\binom kj P(j)$（第 13 题与第 14 题），是
$P(0), \dots, P(k)$ 的整数线性组合。这正是
波利亚对整值多项式的刻画。

\textbf{18.} 设 $Q(X) = P(X + a)$，它是次数 $\leq
n$ 的多项式，且 $Q(0), Q(1), \dots, Q(n) \in \Z$。其关于
$(B_k)_{k \leq n}$ 的坐标为 $c_k = \sum_{j \leq k}(-1)^{k-j}\binom kj
Q(j) \in \Z$（第 14 题只用到 $0, \dots, k
\leq n$）。由第 17 题（$\Leftarrow$）可知 $Q$ 在
整个 $\Z$ 上取整值，故 $P(X) = Q(X - a)$ 亦然。

\textbf{19.} $k$ 个连续整数之积 $m(m-1)
\cdots(m-k+1) = k!\,B_k(m)$，其中 $m \in \Z$，而 $B_k(m) \in
\Z$（第 16 题），故此积可被 $k!$ 整除。

\textbf{20.} $P = \frac{X(X+1)(2X+1)}{6}$ 在 $0,1,2,3$ 处的值为：
$0, 1, 5, 14$。差分表：$\Delta$ 行为 $1, 4, 9$；
$\Delta^2$ 行为 $3, 5$；$\Delta^3$ 行为 $2$。由第 13 题，
\[
P = 0\cdot B_0 + 1\cdot B_1 + 3\,B_2 + 2\,B_3 ,
\]
坐标均为整数：$P$ 为整值多项式（第 17 题），
而其单项式系数 $\frac13, \frac12, \frac16$ 都
不是整数。直接计算：
\begin{align*}
\Delta P &= \frac{(X+1)(X+2)(2X+3) - X(X+1)(2X+1)}{6} \\
&= \frac{(X+1)\bigl[(X+2)(2X+3) - X(2X+1)\bigr]}{6}
= \frac{(X+1)(6X+6)}{6} = (X+1)^2 .
\end{align*}
由逐项相消，$P(m) = \sum_{j=0}^{m-1}\Delta P(j) = \sum_{j=1}^{m}
j^2$（取 $P(0) = 0$），即平方和公式。

\textbf{21.} $2^i$ 在 $i = 0, \dots, n$ 处的值，其差分表
左边缘处恒为 $1$：$\Delta^k$ 作用于数列
$(2^i)$ 后仍为 $(2^i)$（因为 $2^{i+1} - 2^i = 2^i$），故
$\bigl(\Delta^k P\bigr)(0) = 2^0 = 1$（对所有 $k \leq n$），且 $P =
B_0 + B_1 + \dots + B_n$，由第 13 题。于是
\[
P(n+1) = \sum_{k=0}^{n}\binom{n+1}{k}
= 2^{n+1} - \binom{n+1}{n+1} = 2^{n+1} - 1 \neq 2^{n+1}:
\]
这一规律在第一个未被控制的点处即告失效。

\textbf{22.} 由第 12 题，$B_k = \Delta B_{k+1}$，故
\[
\sum_{j=0}^{m-1} B_k(j)
= \sum_{j=0}^{m-1}\bigl(B_{k+1}(j+1) - B_{k+1}(j)\bigr)
= B_{k+1}(m) - B_{k+1}(0) = B_{k+1}(m).
\]
当 $j < k$ 时 $B_k(j)$ 为零，故求和实际上从
$j = k$ 开始：$\sum_{j=k}^{m-1}\binom jk = \binom m{k+1}$，即
曲棍球棒恒等式。

\textbf{23.} 由差分表（或直接展开）可得
\[
X^2 = B_1 + 2 B_2, \qquad X^3 = B_1 + 6 B_2 + 6 B_3
\]
（验算：$B_1 + 2B_2 = X + X(X-1) = X^2$；在 $X = 1, 2, 3$ 处，
第二式给出 $1, 8, 27$）。再由第 22 题得
\[
\sum_{j=0}^{m-1} j^2 = B_2(m) + 2B_3(m)
= \binom m2 + 2\binom m3 = \frac{m(m-1)(2m-1)}{6},
\]
\[
\sum_{j=0}^{m-1} j^3 = B_2(m) + 6B_3(m) + 6B_4(m)
= \binom m2 + 6\binom m3 + 6\binom m4 .
\]
展开最后一式：$\binom m2 + 6\binom m3 + 6\binom
m4 = \frac{m(m-1)}{2}\bigl[1 + 2(m-2) +
\frac{(m-2)(m-3)}{2}\bigr] = \frac{m^2(m-1)^2}{4} = \binom m2^2$.
把 $m$ 换为 $m + 1$：$1^3 + \dots + m^3 =
\bigl(\frac{m(m+1)}2\bigr)^2 = (1 + \dots + m)^2$，此即尼科马库斯
恒等式。

\textbf{24.} 节点取 $0, 1, 2$，多项式为 $X^2$。单项式基
$(1, X, X^2)$ 下的坐标为 $(0, 0, 1)$。拉格朗日基
下的坐标就是取值 $(0, 1, 4)$（第 4 题）。牛顿基
$(1, X, X(X-1))$：差商 $f[0] = 0$，$f[0,1] = 1$，
$f[0,1,2] = \frac{3 - 1}{2} = 1$（第 9 题），故坐标为
$(0, 1, 1)$ --- 事实上 $X + X(X-1) = X^2$。三组基、
三个坐标向量，同一个多项式。

\textbf{25.} (i) 第 4 题自动成立，因为 $(L_i)$ 是一个
\emph{基}：插值的存在性与唯一性恰好
等价于坐标的存在性与唯一性。(ii) 牛顿基
是阶梯形的，故坐标可通过
逐次作商算出 --- 每增加一个新节点只添一项，而不扰动
前面各项 --- 而 $P$ 的拉格朗日坐标即为
取值 $P(x_i)$，无需任何计算。(iii) 两组基均
线性无关，依据
\cref{prop:b1:vspaces:freecriteria}中的两条判别准则：牛顿基看次数两两互异，
拉格朗日基看节点处取值。(iv) 波利亚定理断言
“$P(\Z) \subseteq \Z$”这一关于取值的性质，等价于
基 $(B_k)$ 下坐标的整数性 --- 多项式的
算术性质只有在
与问题相适应的基中才显现。
\end{solution}
